\documentclass[12pt,a4paper]{amsart}

\usepackage{amsmath,amssymb,amsthm}
\usepackage{mathtools}
\usepackage{hyperref}
\usepackage{geometry}
\geometry{margin=2.5cm}

% -------------------------------------------------------
\newtheorem{theorem}{Theorem}[section]
\newtheorem{lemma}[theorem]{Lemma}
\newtheorem{corollary}[theorem]{Corollary}
\newtheorem{proposition}[theorem]{Proposition}
\theoremstyle{definition}
\newtheorem{definition}[theorem]{Definition}
\newtheorem{remark}[theorem]{Remark}

\newcommand{\Z}{\mathbb{Z}}
\DeclareMathOperator{\lcm}{lcm}
\newcommand{\euler}{\varphi}

% -------------------------------------------------------
\begin{document}
% -------------------------------------------------------

\title{On the Incompatibility of Giuga Pseudoprimes and Carmichael Numbers:\\
       A Chinese Remainder Theorem Obstruction}

\author{Kurt Ingwer}
\address{Specialist Mathematics Project, Version Build 43}
\email{specialist-maths@localhost}
\date{\today}

\begin{abstract}
A \emph{Giuga pseudoprime} is a composite integer satisfying a certain
primality criterion due to Giuga~(1950), and a \emph{Carmichael number} is a
composite integer satisfying Fermat's little theorem for every base coprime
to it.  Both classes share the property of being squarefree composites that
``mimic'' primes in specific arithmetic tests.  The question whether a number
can simultaneously belong to both classes — a \emph{Giuga-Carmichael number} —
is open.

We derive a necessary congruence condition for Giuga-Carmichael numbers:
any such~$n$ must satisfy $n \equiv p \pmod{p^2(p-1)}$ for every prime $p \mid n$.
Using the Chinese Remainder Theorem, we show that this system of congruences
is \emph{inconsistent} for every $n$ divisible by all three primes~$3$, $5$,
and~$7$.  As a consequence, no Giuga-Carmichael number can have the set
$\{3, 5, 7\}$ as a subset of its prime factors.

We further prove that all Carmichael numbers are odd, and verify computationally
that no Giuga-Carmichael number exists below $10^4$.
\end{abstract}

\maketitle

% -------------------------------------------------------
\section{Introduction}
% -------------------------------------------------------

\subsection{Background}

Two famous classes of ``prime imposters'' in number theory are:

\begin{definition}
A composite integer $n$ is a \emph{Giuga pseudoprime} if it is squarefree
and satisfies, for every prime $p \mid n$:
\begin{enumerate}
  \item[(W)] $p \mid \dfrac{n}{p} - 1$, \quad and
  \item[(S)] $(p-1) \mid \dfrac{n}{p} - 1$.
\end{enumerate}
\end{definition}

\begin{definition}
A composite integer $n$ is a \emph{Carmichael number} if
$a^{n-1} \equiv 1 \pmod{n}$ for every integer $a$ with $\gcd(a, n) = 1$.
By Korselt's criterion \cite{Korselt1899}, this is equivalent to:
$n$ is squarefree and $(p-1) \mid (n-1)$ for every prime $p \mid n$.
\end{definition}

No Giuga pseudoprime has ever been found.  Carmichael numbers, by contrast,
are known to be infinite in number \cite{Alford1994}.  The question whether
a number can be both — a \emph{Giuga-Carmichael number} — appears to be open.

\subsection{Main results}

\begin{theorem}[Necessary congruence]
\label{thm:congruence}
If $n$ is a Giuga-Carmichael number and $p \mid n$ is prime, then
\[
  n \equiv p \pmod{p^2(p-1)}.
\]
\end{theorem}

\begin{theorem}[CRT obstruction for $\{3, 5, 7\}$]
\label{thm:crt357}
No integer divisible by all three primes $3$, $5$, and $7$ satisfies
the system of congruences in Theorem~\ref{thm:congruence}.
In particular, no Giuga-Carmichael number is divisible by each of $3$, $5$, $7$.
\end{theorem}

% -------------------------------------------------------
\section{Preliminary: Carmichael Numbers are Odd}
\label{sec:odd}
% -------------------------------------------------------

\begin{proposition}
\label{prop:odd}
Every Carmichael number is odd.
\end{proposition}

\begin{proof}
Let $n$ be an even Carmichael number.  Taking $a = -1$ (with $\gcd(-1, n) = 1$):
by definition, $(-1)^{n-1} \equiv 1 \pmod{n}$.  Since $n$ is even, $n - 1$
is odd, so $(-1)^{n-1} = -1$.  Thus $-1 \equiv 1 \pmod{n}$, which gives
$n \mid 2$.  But $n > 2$ is composite — a contradiction.
\end{proof}

By Proposition~\ref{prop:odd}, every Giuga-Carmichael number is odd, so all
its prime factors are odd and at least~$3$.

% -------------------------------------------------------
\section{The Necessary Congruence Condition}
\label{sec:congruence}
% -------------------------------------------------------

\begin{proof}[Proof of Theorem~\ref{thm:congruence}]
Let $n$ be a Giuga-Carmichael number and $p \mid n$ a prime factor.

\medskip
\noindent\textbf{From the Giuga weak condition~(W):}
$p \mid \tfrac{n}{p} - 1$, i.e.\ $\tfrac{n}{p} \equiv 1 \pmod{p}$,
i.e.\ $n \equiv p \pmod{p^2}$.

\medskip
\noindent\textbf{From Korselt's criterion:}
$(p-1) \mid (n-1)$, i.e.\ $n \equiv 1 \pmod{p-1}$.
Since $p = (p-1) + 1 \equiv 1 \pmod{p-1}$, this gives $n \equiv p \pmod{p-1}$.

\medskip
\noindent\textbf{Combining via CRT:}
We have $n \equiv p \pmod{p^2}$ and $n \equiv p \pmod{p-1}$.
Since $\gcd(p^2, p-1) = 1$ (as $p$ is prime: every common divisor~$d$
of $p^2$ and $p-1$ divides $\gcd(p, p-1) = 1$),
the Chinese Remainder Theorem gives
\[
  n \equiv p \pmod{p^2(p-1)}.  \qedhere
\]
\end{proof}

\begin{remark}
The moduli $p^2(p-1)$ for small primes are:
$p=3$: $18$; $p=5$: $100$; $p=7$: $294$; $p=11$: $1210$; $p=13$: $2028$.
\end{remark}

% -------------------------------------------------------
\section{The CRT Obstruction}
\label{sec:crt}
% -------------------------------------------------------

\begin{proof}[Proof of Theorem~\ref{thm:crt357}]
Suppose $n$ is a Giuga-Carmichael number divisible by $3$, $5$, and $7$.
By Theorem~\ref{thm:congruence}:
\begin{align}
  n &\equiv 3 \pmod{18},   \label{eq:p3} \\
  n &\equiv 5 \pmod{100},  \label{eq:p5} \\
  n &\equiv 7 \pmod{294}.  \label{eq:p7}
\end{align}

\noindent\textbf{Step 1: Combining \eqref{eq:p3} and \eqref{eq:p5}.}
We seek $n$ satisfying both.  Write $n = 18k + 3$.  Substituting into
\eqref{eq:p5}:
\[
  18k + 3 \equiv 5 \pmod{100}
  \;\Longrightarrow\;
  18k \equiv 2 \pmod{100}
  \;\Longrightarrow\;
  9k \equiv 1 \pmod{50}.
\]
Since $9 \cdot 39 = 351 = 7 \cdot 50 + 1 \equiv 1 \pmod{50}$, we have
$9^{-1} \equiv 39 \pmod{50}$, giving $k \equiv 39 \pmod{50}$.
Thus $n = 18(50m + 39) + 3 = 900m + 705$, i.e.\
\begin{equation}
  n \equiv 705 \pmod{900}.
  \label{eq:p35}
\end{equation}

\noindent\textbf{Step 2: Checking compatibility with \eqref{eq:p7}.}
The system \eqref{eq:p35}--\eqref{eq:p7} is solvable if and only if
\[
  \gcd(900, 294) \;\Big|\; (7 - 705).
\]
We compute $\gcd(900, 294) = 6$.  And $(7 - 705) = -698 \equiv 4 \pmod{6}$,
since $698 = 116 \cdot 6 + 2$ gives $-698 \equiv -2 \equiv 4 \pmod 6$.

Since $4 \ne 0$, the condition $6 \mid (7 - 705)$ fails.
The system has \emph{no integer solution} — a contradiction.
\end{proof}

\begin{remark}
The contradiction arises from the interplay between the moduli $18$, $100$,
and $294$.  A key observation is that $900 = \lcm(18, 100)$ and $\gcd(900, 294) = 6$,
while $7 - 705 \equiv 4 \pmod{6}$, making the final congruence incompatible.
\end{remark}

% -------------------------------------------------------
\section{Extensions and Open Problems}
\label{sec:open}
% -------------------------------------------------------

\subsection{Further CRT obstructions}

The same method can be applied to other prime triples.  We have verified
computationally:

\begin{proposition}
For the prime triple $\{3, 5, 11\}$: the CRT system is also inconsistent.
For the prime triple $\{5, 7, 13\}$ with smallest factor $5 > 3$: the system
is consistent (a solution $n = 2107105$ exists), showing that the obstruction
is sensitive to the specific choice of prime factors.
\end{proposition}

\subsection{Open problems}

\begin{enumerate}
  \item Is the system of congruences in Theorem~\ref{thm:congruence}
        inconsistent for \emph{every} finite set of odd primes?
  \item Can the structure of Giuga-Carmichael numbers be further constrained
        using strong Giuga conditions~(S) in addition to the weak conditions used
        here?
  \item Is it provable (without computer search) that no Giuga-Carmichael
        number exists below a given explicit bound?
\end{enumerate}

If the answer to~(1) is yes, that would imply the nonexistence of
Giuga-Carmichael numbers unconditionally.

% -------------------------------------------------------
\section{Numerical Verification}
% -------------------------------------------------------

A complete computer search confirms that no integer below $10^4$ is
simultaneously a Giuga pseudoprime and a Carmichael number.  The search
proceeds by:
\begin{enumerate}
  \item Selecting squarefree composite $n$.
  \item Checking the weak Giuga condition for each prime $p \mid n$.
  \item Checking the strong Giuga condition for each prime $p \mid n$.
  \item Checking Korselt's criterion for each prime $p \mid n$.
\end{enumerate}

The known Giuga numbers $\{30, 858, 1722, 66198, \ldots\}$ all fail the strong
Giuga condition~(S) for at least one prime factor, so they are not Giuga
pseudoprimes and a fortiori not Giuga-Carmichael numbers.

The known Carmichael numbers below $10^4$ are $\{561, 1105, 1729, 2465, 2821,
6601, 8911\}$. Direct inspection shows that none of them satisfies even the
weak Giuga condition.

% -------------------------------------------------------
\begin{thebibliography}{10}

\bibitem{Giuga1950}
G.~Giuga,
\emph{Su una presumibile propriet\`a caratteristica dei numeri primi},
Ist.\ Lombardo Sci.\ Lett.\ Rend.\ Cl.\ Sci.\ Mat.\ Nat.\ \textbf{83}
(1950), 511--528.

\bibitem{Korselt1899}
A.~Korselt,
\emph{Probl\`eme chinois},
L'Interm\'ediaire des math\'ematiciens \textbf{6} (1899), 142--143.

\bibitem{Alford1994}
W.~R.~Alford, A.~Granville, and C.~Pomerance,
\emph{There are infinitely many Carmichael numbers},
Ann.\ of Math.\ (2) \textbf{139} (1994), no.\ 3, 703--722.

\bibitem{Borwein1996}
J.~M.~Borwein, P.~B.~Borwein, R.~Girgensohn, and S.~Parnes,
\emph{Giuga's conjecture on primality},
Amer.\ Math.\ Monthly \textbf{103} (1996), no.\ 1, 40--50.

\bibitem{Pomerance1980}
C.~Pomerance, J.~L.~Selfridge, and S.~S.~Wagstaff, Jr.,
\emph{The pseudoprimes to $25 \cdot 10^9$},
Math.\ Comp.\ \textbf{35} (1980), no.\ 151, 1003--1026.

\end{thebibliography}

\end{document}
